\documentclass[tikz,border=5pt]{standalone}
\usetikzlibrary{calc,positioning,intersections,shadows.blur,decorations.pathmorphing,patterns,shapes.callouts}

\definecolor{moonGray}{RGB}{180,180,180}
\definecolor{craterGray}{RGB}{120,120,120}
\definecolor{spaceBlue}{RGB}{10,10,40}
\definecolor{sunlight}{RGB}{255, 215, 0}
\definecolor{earthBlue}{RGB}{60,170,255}
\definecolor{earthGreen}{RGB}{40,140,60}
\definecolor{bedrock}{RGB}{120, 120, 120}
\definecolor{regolith}{RGB}{180, 150, 120}

\begin{document}
\begin{tikzpicture}[scale=1.2]

% White background for the entire area with equal borders
\fill[white] (0,-4.1) rectangle (10.1,3.5);

% Single space box above the surface section
\fill[spaceBlue] (0,0) rectangle (10,3.5);

% Add stars to the space section
\foreach \i in {1,...,75}{
  \fill[white] (rnd*10,rnd*2+1.2) circle (0.02);
}

% Night case surface - extended across the whole frame
\shade[top color=regolith, bottom color=bedrock] (0,-2.5) rectangle (10,0);
\fill[bedrock] (0,-4) rectangle (10,-2.4);

% Add rock fragments randomly distributed across the surface
\pgfmathsetseed{42} % Set seed for reproducible randomness
\foreach \i in {1,...,30} {
    \pgfmathsetmacro{\xrel}{rnd*9}
    \pgfmathsetmacro{\y}{-0.3 + rnd*(-2.2)}
    \pgfmathsetmacro{\size}{0.1 + rnd*0.15}
    
    \fill[bedrock] (\xrel,\y) -- (\xrel+\size,\y-0.05) -- (\xrel+\size+0.05,\y+\size*0.7) -- (\xrel+\size*0.5,\y+\size) -- (\xrel,\y+\size*0.5) -- cycle;
}

% Textured surface line
\draw[black,thick,decorate,decoration={random steps, segment length=3pt, amplitude=0.8pt}] (0,0) -- (10,0);

% Text field for drill core (in situ) measurements - left side
\node[fill=white, draw=black, thick, text width=4cm, align=left, font=\footnotesize\bfseries] at (2.5,-1.7) {
    In Situ Drilling\\
    \vspace{2pt}
    {\normalfont\scriptsize\color{green!60!black}
    • Direct measurement\\
    • Material sampling\\
    • Depth profiling\\
    \vspace{3pt}
    \color{red!70!black}
    • Limited coverage\\
    • Complex equipment\\
    • High mission cost}
};

% Text field for remote sensing - right side
\node[fill=white, draw=black, thick, text width=4cm, align=left, font=\footnotesize\bfseries] at (7.5,-1.7) {
    Remote Sensing\\
    \vspace{2pt}
    {\normalfont\scriptsize\color{green!60!black}
    • Large spatial coverage\\
    • Non-invasive operation\\
    • Rapid data acquisition\\
    \vspace{4pt}
    \color{red!70!black}
    • Model-dependent results\\
    • Limited spatial resolution\\
    • Sources of interference}
};

% Red wavy arrows emitted from surface upwards (thermal radiation)
\foreach \i in {0,...,7} {
    \draw[red,thick,->,decorate,decoration={snake, segment length=6pt, amplitude=2pt}] (1.2 + \i*1.1, 0) -- (1.2 + \i*1.1, 0.4);
}

% Add rover on the left side of the surface
\node[anchor=south] at (2.0,-0.3) {\includegraphics[width=2.0cm]{textures/Rover.png}};

% Add LRO image in top right corner
\node[anchor=north east] at (9.0,3.5) {\includegraphics[width=1.5cm]{textures/LRO.png}};

% Add triangular FOV scan beams from LRO
\draw[yellow, thick, opacity=0.7] (8.4,2.4) -- (7.0,0) -- (7.4,0) -- cycle;
\draw[yellow, thick, opacity=0.7] (8.4,2.4) -- (7.4,0) -- (7.8,0) -- cycle;
\draw[yellow, thick, opacity=0.7] (8.4,2.4) -- (7.8,0) -- (8.2,0) -- cycle;

\end{tikzpicture}
\end{document}